\chapter{Likelihood-Based Learning with Latent Variables}

\section{Maximum Likelihood Estimation}

A central principle in generative modeling is maximum likelihood estimation (MLE). The goal is to choose model parameters that maximize the probability of observed data.

\medskip

\noindent
\textbf{Objective.}  
Given data $\{x_i\}_{i=1}^n$, we maximize:
\[
\max_\theta \; \frac{1}{n} \sum_{i=1}^n \log p_\theta(x_i).
\]

\medskip

\noindent
\textbf{Interpretation.}
\begin{itemize}
    \item Encourages the model to assign high probability to observed data
    \item Equivalent to minimizing the negative log-likelihood
\end{itemize}

\medskip

\noindent
\textbf{Connection to divergence.}  
MLE minimizes the Kullback--Leibler divergence:
\[
\mathrm{KL}(p_{\text{data}} \,\|\, p_\theta).
\]

\medskip

\noindent
\textbf{Gradient.}
\[
\nabla_\theta \log p_\theta(x)
\]
can often be computed efficiently for explicit models.

\medskip

\noindent
\textbf{Insight.}  
MLE provides a principled way to learn probabilistic models from data.

\section{Hidden Variables and Marginalization}

Latent variable models introduce hidden variables $z$ to capture underlying structure.

\medskip

\noindent
\textbf{Marginal likelihood.}
\[
p_\theta(x) = \int p_\theta(x, z) \, dz = \int p_\theta(x \mid z) p(z) \, dz.
\]

\medskip

\noindent
\textbf{Log-likelihood.}
\[
\log p_\theta(x) = \log \int p_\theta(x, z) \, dz.
\]

\medskip

\noindent
\textbf{Interpretation.}
\begin{itemize}
    \item $z$ explains variability in $x$
    \item The model integrates over all possible latent explanations
\end{itemize}

\medskip

\noindent
\textbf{Challenge.}  
The integral over $z$ is often intractable, especially in high dimensions.

\medskip

\noindent
\textbf{Posterior distribution.}
\[
p_\theta(z \mid x) = \frac{p_\theta(x, z)}{p_\theta(x)}.
\]

\medskip

\noindent
\textbf{Insight.}  
Learning requires both evaluating the marginal likelihood and reasoning about latent variables.

\section{Intractability and Approximation}

Exact computation of the marginal likelihood is typically infeasible.

\medskip

\noindent
\textbf{Sources of intractability.}
\begin{itemize}
    \item High-dimensional latent space
    \item Complex likelihood functions
\end{itemize}

\medskip

\noindent
\textbf{Approximation strategies.}
\begin{itemize}
    \item Sampling-based methods (e.g., Monte Carlo)
    \item Deterministic approximations
    \item Variational methods
\end{itemize}

\medskip

\noindent
\textbf{Monte Carlo approximation.}
\[
p_\theta(x) \approx \frac{1}{N} \sum_{i=1}^N p_\theta(x \mid z_i), \quad z_i \sim p(z).
\]

\medskip

\noindent
\textbf{Limitations.}
\begin{itemize}
    \item High variance estimates
    \item Computationally expensive
\end{itemize}

\medskip

\noindent
\textbf{Alternative viewpoint.}  
Rather than approximating the integral directly, we approximate the posterior $p_\theta(z \mid x)$.

\medskip

\noindent
\textbf{Insight.}  
Approximation is unavoidable in latent variable models and shapes both theory and practice.

\section{Variational Principles}

Variational methods provide a powerful framework for approximating intractable quantities.

\medskip

\noindent
\textbf{Variational distribution.}  
Introduce a tractable distribution $q_\phi(z \mid x)$ to approximate the true posterior.

\medskip

\noindent
\textbf{Key identity.}
\[
\log p_\theta(x)
=
\mathbb{E}_{q_\phi(z \mid x)}\left[\log \frac{p_\theta(x, z)}{q_\phi(z \mid x)}\right]
+
\mathrm{KL}(q_\phi(z \mid x) \,\|\, p_\theta(z \mid x)).
\]

\medskip

\noindent
\textbf{Evidence lower bound (ELBO).}
\[
\mathcal{L}(\theta, \phi)
=
\mathbb{E}_{q_\phi(z \mid x)}[\log p_\theta(x, z) - \log q_\phi(z \mid x)].
\]

\medskip

\noindent
\textbf{Lower bound property.}
\[
\mathcal{L}(\theta, \phi) \leq \log p_\theta(x).
\]

\medskip

\noindent
\textbf{Decomposition.}
\[
\mathcal{L}(\theta, \phi)
=
\mathbb{E}_{q_\phi(z \mid x)}[\log p_\theta(x \mid z)]
-
\mathrm{KL}(q_\phi(z \mid x) \,\|\, p(z)).
\]

\medskip

\noindent
\textbf{Interpretation.}
\begin{itemize}
    \item First term: reconstruction quality
    \item Second term: regularization toward the prior
\end{itemize}

\medskip

\noindent
\textbf{Optimization.}
\begin{itemize}
    \item Maximize ELBO with respect to $\theta$ and $\phi$
    \item Jointly learn generative and inference models
\end{itemize}

\medskip

\noindent
\textbf{Insight.}  
Variational methods transform intractable inference into an optimization problem.

\section{A Unifying Perspective}

Likelihood-based learning with latent variables involves several interconnected ideas:

\begin{itemize}
    \item \textbf{MLE:} learning from observed data
    \item \textbf{Latent variables:} modeling hidden structure
    \item \textbf{Approximation:} handling intractable integrals
    \item \textbf{Variational inference:} optimization-based approximation
\end{itemize}

\medskip

\noindent
\textbf{Core tradeoff.}
\[
\text{Accuracy of approximation} \quad \text{vs} \quad \text{computational tractability}.
\]

\medskip

\noindent
\textbf{Key insight.}  
Introducing latent variables increases modeling power but requires approximate inference techniques.

\section{Outlook}

This chapter introduced likelihood-based learning with latent variables:
\begin{itemize}
    \item maximum likelihood estimation,
    \item marginalization over hidden variables,
    \item challenges of intractability,
    \item variational principles and the ELBO.
\end{itemize}

\medskip

In the next chapter, we apply these ideas to variational autoencoders, a practical and widely used class of generative models.

\medskip

\noindent
\textbf{Guiding principle.}  
Complex probabilistic models require approximate inference, and variational methods provide a principled and scalable solution.